%%%%%%%%%%%%%%%%%%%%%%%%%%%%%%%%%%%%%%%%%%%%%%%%%%%%%%%%%%%%%%%%%%%%%%
% Lớp 8 | Chương III. Tứ giác | Hình thoi và hình vuông
% Dạng 8H3?14-2: Tính chất, dấu hiệu nhận biết hình vuông
%%%%%%%%%%%%%%%%%%%%%%%%%%%%%%%%%%%%%%%%%%%%%%%%%%%%%%%%%%%%%%%%%%%%%%


% nguồn: De_AI_Toan8_cuoi_hoc_ki_i_1 (nhập 21/11/2026)
\begin{bt}%[8H3H14-2]%[ma:6771e4d7]
	Cho hình vuông $ABCD$ có độ dài cạnh bằng $6\text{ cm}$.
a) Tính chu vi và diện tích của hình vuông $ABCD$.
b) Tính độ dài đường chéo $AC$ và số đo góc $BAC$.
	\loigiai{
		a) Chu vi hình vuông $ABCD$ là $4 \cdot 6 = 24\text{ cm}$. \\
Diện tích hình vuông $ABCD$ là $6^2 = 36\text{ cm}^2$. \\
b) Áp dụng định lí Pythagore trong tam giác $ABC$ vuông tại $B$: $AC^2 = AB^2 + BC^2 = 6^2 + 6^2 = 72 \Rightarrow AC = 6\sqrt{2}\text{ cm}$. \\
Vì $ABCD$ là hình vuông nên đường chéo $AC$ là tia phân giác của góc $BAD = 90^\circ$, suy ra $\widehat{BAC} = \dfrac{90^\circ}{2} = 45^\circ$.
	}
\end{bt}

% nguồn: De_AI_Toan8_cuoi_hoc_ki_i_1 (nhập 21/11/2026)
\begin{bt}%[8H3H14-2]%[ma:b2a8b560]
	Cho tam giác $ABC$ vuông tại $A$, đường phân giác của góc $A$ cắt cạnh $BC$ tại $D$. Kẻ $DE \perp AB$ tại $E$ và $DF \perp AC$ tại $F$. Chứng minh tứ giác $AEDF$ là hình vuông.
	\loigiai{
		Xét tứ giác $AEDF$ có $\widehat{EAF} = 90^\circ$ (tam giác $ABC$ vuông tại $A$), $\widehat{AED} = 90^\circ$ ($DE \perp AB$), $\widehat{AFD} = 90^\circ$ ($DF \perp AC$). \\
Tứ giác $AEDF$ có ba góc vuông nên $AEDF$ là hình chữ nhật. \\
Mặt khác, $AD$ là đường phân giác của góc $EAF$. \\
Hình chữ nhật $AEDF$ có một đường chéo là đường phân giác của một góc nên $AEDF$ là hình vuông.
	}
\end{bt}

% nguồn: De_AI_Toan8_cuoi_hoc_ki_i_1 (nhập 21/11/2026)
\begin{bt}%[8H3H14-2]%[ma:e9eecb94]
	Cho hình vuông $ABCD$. Trên các cạnh $AB$, $BC$, $CD$, $DA$ lần lượt lấy các điểm $E$, $F$, $G$, $H$ sao cho $AE = BF = CG = DH$.
a) Chứng minh bốn tam giác vuông $AEH$, $BFE$, $CGF$, $DHG$ bằng nhau.
b) Chứng minh tứ giác $EFGH$ là hình vuông.
	\loigiai{
		a) Vì $ABCD$ là hình vuông nên $AB = BC = CD = DA$ và $\widehat{A} = \widehat{B} = \widehat{C} = \widehat{D} = 90^\circ$. \\
Mà $AE = BF = CG = DH$ nên các đoạn còn lại cũng bằng nhau: $EB = FC = GD = HA$. \\
Suy ra bốn tam giác vuông $AEH$, $BFE$, $CGF$, $DHG$ bằng nhau theo trường hợp hai cạnh góc vuông (c.g.c). \\
b) Từ bốn tam giác bằng nhau ở câu a) suy ra $EH = FE = GF = HG$, nên $EFGH$ là hình thoi. \\
Lại có $\widehat{AEH} + \widehat{BEF} = \widehat{AEH} + \widehat{AHE} = 90^\circ$ (do $\widehat{BEF} = \widehat{AHE}$). \\
Suy ra $\widehat{HEF} = 180^\circ - (\widehat{AEH} + \widehat{BEF}) = 180^\circ - 90^\circ = 90^\circ$. \\
Hình thoi $EFGH$ có một góc vuông nên là hình vuông.
	}
\end{bt}
